\section{Circuit Complexity --- Logspace Uniform Adj Basic}
%%% generated by the blueprint pipeline from TCSlib.Complexity.CircuitComplexity.LogspaceUniformAdjBasic
\subsection*{Circuit Complexity --- Logspace Uniform Adj Basic: Introduction}

Declarations of the Lean module \texttt{TCSlib.Complexity.CircuitComplexity.LogspaceUniformAdjBasic}.

\subsection*{Circuit Complexity --- Logspace Uniform Adj Basic: Results}

\begin{definition}[Vertex type]
\lean{BoolCircuit.DAGCircuit.vtype}
\label{BoolCircuit.DAGCircuit.vtype}
\leanok
\uses{BoolCircuit.GateKind}
For a circuit with $n$ input vertices and a list of gates, define the type of vertex $i$
to be absent when $i<n$. When $i\ge n$, it is the kind of the gate at list position
$i-n$, if that position exists, and absent otherwise.
\end{definition}

\begin{definition}[Circuit edge relation]
\lean{BoolCircuit.DAGCircuit.Edge}
\label{BoolCircuit.DAGCircuit.Edge}
\leanok
For a circuit with $n$ input vertices, there is an edge from vertex $i$ to vertex $j$
exactly when $j\ge n$ and the gate at list position $j-n$ exists and has $i$ among its
argument indices.
\end{definition}

\begin{definition}[Canonical circuit]
\lean{BoolCircuit.DAGCircuit.IsCanonical}
\label{BoolCircuit.DAGCircuit.IsCanonical}
\leanok
\uses{BoolCircuit.DAGCircuit.size}
A circuit is canonical if every gate’s argument indices are in strictly increasing order
and its output index plus one equals its size, where the size is the number of input
vertices plus the number of gates.
\end{definition}

\begin{lemma}[Expanded circuit encoding]
\lean{BoolCircuit.DAGCircuit.encode_eq}
\label{BoolCircuit.DAGCircuit.encode_eq}
\leanok
\uses{BoolCircuit.DAGCircuit.encode, BoolCircuit.DAGGate.encode, BoolCircuit.encodeList, BoolCircuit.encodeNat}
\difficulty{1}
For a directed acyclic Boolean circuit $C$ on $n$ input vertices with output vertex $o$,
its Boolean-string encoding is the concatenation of $n$ true bits followed by a false
bit, the encoded gate list, and $o$ true bits followed by a false bit.
\end{lemma}

\begin{definition}[Circuit size language]
\lean{BoolCircuit.DAGCircuitFamily.sizeLang}
\label{BoolCircuit.DAGCircuitFamily.sizeLang}
\leanok
\uses{BoolCircuit.DAGCircuit.size, Turing.pairEncode}
The size language of a circuit family consists of the strings that encode a pair whose
first component is a unary string of $n$ true symbols and whose second component is the
binary representation of a natural number $i$, for some $n$ and $i$, with $i$ strictly
less than the size of the family’s circuit at parameter $n$.
\end{definition}

\begin{definition}[Gate type language]
\lean{BoolCircuit.DAGCircuitFamily.typeLang}
\label{BoolCircuit.DAGCircuitFamily.typeLang}
\leanok
\uses{BoolCircuit.DAGCircuit.size, BoolCircuit.DAGCircuit.vtype, BoolCircuit.GateKind, Turing.pairEncode}
For a circuit family and a gate type $t$ (or absence of a gate type), the corresponding
language consists of strings encoding a pair of a unary string of $n$ true symbols and
the binary representation of a natural number $i$, for some $n$ and $i$, such that $i$
is strictly less than the circuit size at $n$ and vertex $i$ has type $t$ in that
circuit.
\end{definition}

\begin{definition}[Circuit edge language]
\lean{BoolCircuit.DAGCircuitFamily.edgeLang}
\label{BoolCircuit.DAGCircuitFamily.edgeLang}
\leanok
\uses{BoolCircuit.DAGCircuit.Edge, Turing.pairEncode}
The edge language of a circuit family consists of strings encoding a unary string of $n$
true symbols together with the pair of binary representations of natural numbers $i$ and
$j$, for some $n,i,j$, such that there is an edge from vertex $i$ to vertex $j$ in the
family’s circuit at parameter $n$.
\end{definition}

\begin{definition}[Logspace adjacency data]
\lean{BoolCircuit.DAGCircuitFamily.HasLogspaceAdjacency}
\label{BoolCircuit.DAGCircuitFamily.HasLogspaceAdjacency}
\leanok
\uses{BoolCircuit.DAGCircuitFamily.IsPolySize, BoolCircuit.DAGCircuitFamily.edgeLang, BoolCircuit.DAGCircuitFamily.sizeLang, BoolCircuit.DAGCircuitFamily.typeLang, Complexity.LOGSPACE}
A circuit family has logspace adjacency data if it has polynomial size, its size
language is decidable in logspace, each of its gate type languages is decidable in
logspace, and its edge language is decidable in logspace.
\end{definition}

\begin{lemma}[Membership in the size language]
\lean{BoolCircuit.DAGCircuitFamily.mem_sizeLang}
\label{BoolCircuit.DAGCircuitFamily.mem_sizeLang}
\leanok
\uses{BoolCircuit.DAGCircuit.size, BoolCircuit.DAGCircuitFamily.sizeLang, Complexity.LogProg.bits_injective, Turing.pairEncode, Turing.pairEncode_replicate_inj}
\difficulty{4}
For a family of Boolean circuits, a natural number $i$ satisfies $i<\mathrm{size}(C_n)$
if and only if the string encoding $n$ copies of the bit $1$ followed by the binary
representation of $i$ belongs to the family’s size language.
\end{lemma}

\begin{lemma}[Membership in the gate type language]
\lean{BoolCircuit.DAGCircuitFamily.mem_typeLang}
\label{BoolCircuit.DAGCircuitFamily.mem_typeLang}
\leanok
\uses{BoolCircuit.DAGCircuit.size, BoolCircuit.DAGCircuit.vtype, BoolCircuit.DAGCircuitFamily.typeLang, BoolCircuit.GateKind, Complexity.LogProg.bits_injective, Turing.pairEncode, Turing.pairEncode_replicate_inj}
\difficulty{4}
For a family of Boolean circuits, a gate type $t$ (or the absence of a gate type) and
natural numbers $n,i$, the encoding of $n$ copies of $1$ followed by the binary
representation of $i$ belongs to the type language for $t$ if and only if
$i<\mathrm{size}(C_n)$ and vertex $i$ of $C_n$ has type $t$.
\end{lemma}

\begin{lemma}[Membership in the edge language]
\lean{BoolCircuit.DAGCircuitFamily.mem_edgeLang}
\label{BoolCircuit.DAGCircuitFamily.mem_edgeLang}
\leanok
\uses{BoolCircuit.DAGCircuit.Edge, BoolCircuit.DAGCircuitFamily.edgeLang, Complexity.LogProg.bits_injective, Turing.pairEncode, Turing.pairEncode_injective, Turing.pairEncode_replicate_inj}
\difficulty{4}
Let $C$ be a family of Boolean circuits, and let $n,i,j$ be natural numbers. The
encoding of the unary string of length $n$ paired with the encodings of $i$ and $j$
belongs to the edge language of $C$ if and only if vertex $i$ is an argument of the gate
at vertex $j$ in the $n$-input circuit, with $j$ a gate vertex.
\end{lemma}

\begin{definition}[Canonical circuit family]
\lean{BoolCircuit.DAGCircuitFamily.IsCanonical}
\label{BoolCircuit.DAGCircuitFamily.IsCanonical}
\leanok
\uses{BoolCircuit.DAGCircuit.IsCanonical}
A circuit family is canonical if, for every natural number $n$, its circuit at parameter
$n$ is canonical: each gate’s argument indices are strictly increasing, and the output
index plus one equals the circuit’s size, defined as $n$ plus the number of gates.
\end{definition}
